% everything is in seperate documentClass
\documentclass{./tpl/styles/dhbook}                                         	

%%%%%%%%%%%%%%%%%%%%%%% Definitionen bzgl. der Arbeit %%%%%%%%%%%%%%%%%%%%%%%
\dhlogo{./img/dhbw_logo.png}
\type{1. Projektarbeit}           
\title{MeinThema}
\subtitle{Themenuntertitel}
\author{Max Mustermann}
\kurs{Muster11}
\studiengang{Informationstechnik}
\fspo{Bachelor of Engenering, Praxismodul I}
\date{\today}
\schoolSupervisor{Prof. Schnattinger}
\companyName{Muster GmbH}
\companyAddress{Musterhausen}
\companyContact{MyCompany\\
		MeineStrasse 23\\
		87654 NeuHausen\\
		Tel: 05421/4324-22\\
		Web: \url{mycompany.de}}
\companyLogo{./img/mycompany_logo.png}
\companySupervisor{Prof. Dr. Donald Ervin Knuth}

\issperrvermerk


%%%%%%%%%%%%% title, Index, Abkürzungsverzeichnis und Glossar erstellen %%%%%%%%%%%%



\makeindex
\makeglossaries

%\newglossaryentry{<label>}{<settings>}

\newglossaryentry{computer}
{
  name=computer,
  description={is a programmable machine that receives input,
               stores and manipulates data, and provides
               output in a useful format}
}

\newglossaryentry{LaTeX}
{
  name=\LaTeX,
  description={is a document markup language and document preparation system for the TeX typesetting program. Within the typesetting system, its name is styled as LaTeX. The term LaTeX refers only to the language in which documents are written, not to the editor used to write those documents. }
}


%%%%%%%%%%%%%%%%%%%%%%%%%%%%%%% PDF-Optionen %%%%%%%%%%%%%%%%%%%%%%%%%%%%%%%%
\hypersetup{
	bookmarksopen=false,
	bookmarksnumbered=true,
	bookmarksopenlevel=0,
	pdftitle=\valueTitle,
	pdfsubject=\valueSubtitle,
	pdfauthor=\valueAuthor,
	pdfborder=0,
	pdfcreator=\valueAuthor,
	pdfstartview=Fit,
	pdfpagelayout=SinglePage
}

\begin{document}


	% Aktiviere blankes Seitenlayout
    \pagestyle{fancy}	   
    % Roemische Seitenzahlen
    \pagenumbering{Roman}
    
    % erzeuge Titelseite + ggf Sperrvermerk
    \maketitle
    
	
    % Zusammenfassungen
    

\phantomsection
\addcontentsline{toc}{chapter}{Zusammenfassung}%
\chapter*{Zusammenfassung}%
\label{sec:zusammenfassung-ger}%
\pdfbookmark[1]{Deutsch}{sec:zusammenfassung-ger}%
%\addcontentsline{toc}{subsection}{Deutsch}

 Lorem ipsum dolor sit amet, consectetur adipiscing elit. Morbi sed diam mauris. Aenean eleifend ultricies nibh sit amet pharetra. Aenean non tellus libero, commodo sodales quam. Curabitur id nunc quis massa euismod lacinia. Nunc eu nibh augue, at suscipit nisl. Cras imperdiet, massa ornare volutpat dapibus, nisi risus adipiscing quam, vel pharetra metus leo id risus. Fusce sit amet quam odio. Vivamus ultrices egestas erat, sit amet pellentesque nisl feugiat a. Morbi tempor lorem ac tellus auctor euismod.  

\par

\newpage  

\chapter*{Summary}%
\label{sec:zusammenfassung-eng}%
\pdfbookmark[1]{Englisch}{sec:zusammenfassung-eng}%
 Lorem ipsum dolor sit amet, consectetur adipiscing elit. Morbi sed diam mauris. Aenean eleifend ultricies nibh sit amet pharetra. Aenean non tellus libero, commodo sodales quam. Curabitur id nunc quis massa euismod lacinia. Nunc eu nibh augue, at suscipit nisl. Cras imperdiet, massa ornare volutpat dapibus, nisi risus adipiscing quam, vel pharetra metus leo id risus. Fusce sit amet quam odio. Vivamus ultrices egestas erat, sit amet pellentesque nisl feugiat a. Morbi tempor lorem ac tellus auctor euismod.  
\par 


    % einfuegen des Vorwortes
    
\phantomsection
\chapter*{Vorwort}%
\addcontentsline{toc}{chapter}{Vorwort}%
Ganz vielen Dank an alle! Und vor allem an die, die mich unterstützt haben.

    
    % erzeuge Inhaltsverzeichnis
    \tableofcontents
    \newpage

    % erzeige Illustartionsverzeichnis (Abbildungen, tabellen Quellcode)
    \listofillustrations

    %% erzeuge Abkuerzungsverzeichnis
    \listglossaries

    % Aktiviere FancyHeader
    \pagestyle{fancy}	   
    % Arabische Seitenzahlen
    \pagenumbering{arabic}	

    % erstes Kapitel
    \chapter{{\LaTeX}-Vorlage}
Das  hier ist eine kleine {\LaTeX}-Vorlage fuer die DHBW Loerach.
Speziell fuer den fachbereich Informatik.

\section{Danke Dounald}
Dank Dounald Knuth\cite{knuth} der Erfinder von \TeX~ koennen wir hier ganz schoen, tolle und Fancy Ausarbeitungen schreiben ohne einen gedanken an das Formatieren zu verschweden.


\newpage
\section{Aufbau der Vorlage}
Die Vorlage ist von der Standart Vorlage book.cls abgeleitet und ist direkt auf DIN A4 gesetzt.
Auch die Schriftart ist auf 12pt und 1,5 Zeilebabstand gesetzt.

\section{Beispiele}
Nun sollen einzelne Beispiele folgen, wie Bilder, Tabelle, Auflistungen, Programmcode, Zitate 

\subsection{Bilder}
ein Bild wird ganz einfach eingebunden, darum herum wird ein \textit{figure}  gesetzt damit die Abbildung auf auf der Abbildungsauflistung erscheint.\\
Mit \textit{caption} wird eine Abbildungsunterschrift angegeben und mit \textit{label} wird ein label vergeb mit dem man auf eine Abbildung verweisen kann.
\begin{figure}[hb] \centering
  \includegraphics[width=3cm]{./img/tux.png}
  \caption{eine Bild Unterschrift}
  \label{tux_1}
\end{figure}
Als Beispiel eines label kann in \ref{tux_1} gesehen werden das dort ein Pinguin abgebildet wird.


\subsection{Tabelle}
Eine Tabelle wird wie bei einer Abbildung wie folgt dargestellt: \\

\begin{table}[hb] \centering
\begin{tabular}{|r|r|}
\hline
Zahl & Quatsch \\
\hline
1  &  0  \\
\hline
2  &  1  \\
\hline
3  &  2  \\
\hline
4  &  3  \\
\hline

\end{tabular}

\caption{eine QuatschTabelle}
\label{tabquatsch}
\end{table}



\section{Glossar}
Das Glossar listet alle Verwendeten Begriffe auf und verlinkt diese auch im Text. Da ist \gls{LaTeX} echt cool und es macht spass mit einem \gls{computer} zuarbeiten.




    % zweites Kapitel
    
\chapter{First}
i\index{Lorem ipsum} \cite{knuth}dolor sit amet, consectetur adipiscing elit. Nullam eget est id dolor ullamcorper dictum. Praesent ornare dictum diam vitae mollis. Aliquam tristique nibh in ante pellentesque non sollicitudin massa aliquam. Integer condimentum urna quis est imperdiet sit amet egestas nisi 
\section{Kuchen}
sollicitudin. Ut et eros nec orci consequat varius. Quisque augue nibh, pulvinar nec adipiscing id, dapibus ut mi. Sed at elementum urna. Mauris id leo et sem condimentum commodo. Ut ligula nunc, aliquet id aliquet ac, pellentesque at turpis. Mauris nibh sem, convallis nec auctor ac, tincidunt quis dolor. Sed leo elit, lacinia vel pharetra tempus, laoreet id sem. Donec at massa eros.

\begin{lstlisting}[caption={Week},label=week1]
if(weekend){
  for(int i; i < 10 ; i++) {
    printf("Hello Weekend\n");
  }
else{
    printf("I want a weekend\n!");
}	
\caption{Weekend3}
\label{week3}
\end{lstlisting}

\index{sollicitudin}. Ut et eros nec orci consequat varius. Quisque augue nibh, pulvinar nec adipiscing id, dapibus ut mi. Sed at elementum urna. Mauris id leo et sem condimentum commodo. Ut ligula nunc, aliquet id aliquet ac, pellentesque at turpis. Mauris nibh sem, convallis nec auctor ac, tincidunt quis dolor. Sed leo elit, lacinia vel pharetra tempus, laoreet id sem. Donec at massa eros.

\begin{enumerate}
  \item The first item
  \item The second item
  \item The third etc \ldots
\end{enumerate}

\section{Schokolade}
Ut vehicula \index{porta} laoreet. Donec sagittis erat nec velit volutpat laoreet. Pellentesque condimentum leo vitae augue pellentesque lacinia. Integer tincidunt suscipit rutrum. Pellentesque volutpat orci vel ligula egestas condimentum. Mauris sit amet libero lacus. In sit amet metus eget metus euismod pharetra. Nulla nisi lacus, tempor ut hendrerit sed, adipiscing id felis. Integer vulputate lectus quis felis sodales ut placerat lacus fringilla.
\begin{lstlisting}[caption={Useless code},label=useless]
if(weekend){
  for(int i; i < 10 ; i++) {
    printf("Hello Weekend\n");
  }
else{
    printf("I want a weekend\n!");
}	
%\caption{Weekend2}
%\label{week2}
\end{lstlisting}

Curabitur ac sapien ante. Phasellus ac diam ut lacus semper scelerisque quis at felis. Donec varius orci a dui fringilla vitae suscipit massa placerat. Aliquam interdum dui vel neque aliquam nec euismod eros placerat. Morbi id nibh erat. Class aptent taciti sociosqu ad litora torquent per conubia nostra, per inceptos himenaeos. In hac habitasse platea dictumst. Praesent velit augue, 
\subsection{Lindt}
Mauris dignissim lectus id tortor tempus sed ultricies nibh dignissim. Praesent id eleifend risus. In hac habitasse platea dictumst. Morbi ac nibh urna, eu pharetra nunc. Lorem ipsum dolor sit amet, consectetur adipiscing elit. Curabitur egestas, odio vel scelerisque vulputate, nibh nibh scelerisque risus, sit amet congue tellus metus sed libero. Praesent fringilla justo at lectus auctor ac euismod sapien lobortis. Nullam sed convallis leo.
\begin{table} \centering
\begin{tabular}{|r|r|}
\hline
Bit-Muster & Zeichen \\
\hline
0000  &  0  \\
\hline
0001  &  1  \\
\hline
0010  &  2  \\
\hline
0011  &  3  \\
\hline
0100  &  4  \\
\hline
0101  &  5  \\
\hline
0110  &  6  \\
\hline
0111  &  7  \\
\hline
\end{tabular}

\caption{eine Tabelle}
\label{tabLindt}
\end{table}

\subsection{Milka}
sollicitudin. Ut et eros \cite{knuth}nec orci consequat varius. Quisque augue nibh, pulvinar nec adipiscing id, dapibus ut mi. Sed at elementum urna. Mauris id leo et sem condimentum commodo. Ut ligula nunc,
\begin{figure} \centering
  \includegraphics[width=6cm]{./img/tux.png}
  \caption{eine Bild Unterschrift}
  \label{tux1}
\end{figure}


 aliquet id aliquet ac, pellentesque at turpis. Mauris nibh sem, convallis nec auctor ac, tincidunt quis dolor. Sed leo elit, lacinia vel pharetra tempus, laoreet id sem. Donec at massa eros.
\begin{table} \centering
\begin{tabular}{|r|r|}
\hline
Bit-Muster & Zeichen \\
\hline
0000  &  0  \\
\hline
0001  &  1  \\
\hline
0010  &  2  \\
\hline
0011  &  3  \\
\hline
0100  &  4  \\
\hline
0101  &  5  \\
\hline
0110  &  6  \\
\hline
0111  &  7  \\
\hline
\end{tabular}
\begin{tabular}{|r|r|}
\hline
Bit-Muster & Zeichen \\
\hline
1000  &  8  \\
\hline
1001  &  9  \\
\hline
1010  &  A  \\
\hline
1011  &  B  \\
\hline
1100  &  C  \\
\hline
1101  &  D  \\
\hline
1110  &  E  \\
\hline
1111  &  F  \\
\hline
\end{tabular}
\caption{Binärzahl zu Hexadezimalzahl}
\label{hex}
\end{table}

Curabitur ac sapien ante. Phasellus ac diam ut lacus semper scelerisque quis at felis. Donec varius orci a dui fringilla vitae suscipit massa placerat. Aliquam interdum dui vel neque aliquam nec euismod eros placerat. Morbi id nibh erat. Class aptent taciti sociosqu ad litora torquent per conubia nostra, per inceptos himenaeos. In hac habitasse platea dictumst. Praesent velit augue, adipiscing dapibus faucibus in, accumsan a massa.
\subsubsection{Loerrach}
Phasellus cursus, eros et scelerisque euismod, enim arcu rhoncus neque, luctus accumsan odio est nec neque. Quisque ultricies mi ultricies risus adipiscing sit amet suscipit odio aliquet. Curabitur 

\begin{lstlisting}[caption={week3},label=weekend34]
if(weekend){
  for(int i; i < 10 ; i++) {
    printf("Hello Weekend\n");
  }
else{
    printf("I want a weekend\n!");
}	
\caption{Weekend1}
\label{week1}
\end{lstlisting}

aliquam euismod mauris. Ut imperdiet lacinia sapien, id condimentum ligula imperdiet eu. Vivamus mollis posuere quam, sed bibendum massa aliquet sed. Donec volutpat luctus mi at commodo. Vivamus faucibus elementum sem a posuere. Donec feugiat, elit in dapibus pulvinar, nibh mi faucibus ipsum, ac volutpat lectus justo nec tellus. Cras in suscipit erat. Vivamus quis dui sem. Nunc posuere ullamcorper tellus, sodales vestibulum massa eleifend quis. Vivamus quis sapien nec metus convallis vulputate quis ac ligula. 

\begin{table} \centering
\begin{tabular}{|r|r|}
\hline
Bit-Muster & Zeichen \\
\hline
0000  &  0  \\
\hline
0001  &  1  \\
\hline
0010  &  2  \\
\hline
0011  &  3  \\
\hline
0100  &  4  \\
\hline
0101  &  5  \\
\hline
0110  &  6  \\
\hline
0111  &  7  \\
\hline
\end{tabular}
\caption{Noch eine Tabelle}
\label{tabloe}
\end{table}


Curabitur ac sapien ante. Phasellus ac diam ut lacus semper scelerisque quis at felis. Donec varius orci a dui fringilla vitae suscipit massa placerat. Aliquam interdum dui vel neque aliquam nec euismod eros placerat. Morbi id nibh erat. Class aptent taciti sociosqu ad litora torquent per conubia nostra, per inceptos himenaeos. In hac habitasse platea dictumst. Praesent velit augue, adipiscing dapibus faucibus in, accumsan a massa.

\begin{figure} \centering
  \includegraphics[width=6cm]{./img/tux.png}
  \caption{eine Bild Unterschrift}
  \label{tux2}
\end{figure}

\section{Kaffee}
Curabitur ac sapien ante. Phasellus ac diam ut lacus semper scelerisque quis at felis. Donec varius orci a dui fringilla vitae suscipit massa placerat. Aliquam interdum dui vel neque aliquam nec euismod eros placerat. Morbi id nibh erat. Class aptent taciti sociosqu ad litora torquent per conubia nostra, per inceptos himenaeos. In hac habitasse platea dictumst. Praesent velit augue, adipiscing dapibus faucibus in, accumsan a massa.
\begin{figure} \centering
  \includegraphics[width=6cm]{./img/tux.png}
  \caption{eine Bild Unterschrift}
  \label{tux3}
\end{figure}


\begin{itemize}

\item Here is a bullet.

\item Here is another bullet.

\end{itemize}


\subsection{Java} Curabitur ac sapien ante. Phasellus ac diam ut lacus semper scelerisque quis at felis. Donec varius orci a dui fringilla vitae suscipit massa placerat. Aliquam interdum dui vel neque aliquam nec euismod eros placerat. Morbi id nibh erat. Class aptent taciti sociosqu ad litora torquent per conubia nostra, per inceptos himenaeos. In hac habitasse platea dictumst. Praesent velit augue, adipiscing dapibus faucibus in, accumsan a massa.
Curabitur ac sapien ante. Phasellus ac diam ut lacus semper scelerisque quis at felis. Donec varius orci a dui fringilla vitae suscipit massa placerat. Aliquam interdum dui vel neque aliquam nec euismod eros placerat. Morbi id nibh erat. Class aptent taciti sociosqu ad litora torquent per conubia nostra, per inceptos himenaeos. In hac habitasse platea dictumst. Praesent velit augue, adipiscing dapibus faucibus in, accumsan a massa.



\chapter{Second}
iLorem ipsum dolor sit amet, consectetur adipiscing elit. Nullam eget est id dolor ullamcorper dictum. Praesent ornare dictum diam vitae mollis. Aliquam tristique nibh in ante pellentesque non sollicitudin massa aliquam. Integer condimentum urna quis est imperdiet sit amet egestas nisi sollicitudin. Ut et eros nec orci consequat varius. Quisque augue nibh, pulvinar nec adipiscing id, dapibus ut mi. Sed at elementum urna. Mauris id leo et sem condimentum commodo. Ut ligula nunc, aliquet id aliquet ac, pellentesque at turpis. Mauris nibh sem, convallis nec auctor ac, tincidunt quis dolor. Sed leo elit, lacinia vel pharetra tempus, laoreet id sem. Donec at massa eros.

\section{Schokolade}
Ut vehicula porta laoreet. Donec sagittis erat nec velit volutpat laoreet. Pellentesque condimentum leo vitae augue pellentesque lacinia. Integer tincidunt suscipit rutrum. Pellentesque volutpat orci vel ligula egestas condimentum. Mauris sit amet libero lacus. In sit amet metus eget metus euismod pharetra. Nulla nisi lacus, tempor ut hendrerit sed, adipiscing id felis. Integer vulputate lectus quis felis sodales ut placerat lacus fringilla.
Curabitur ac sapien ante. Phasellus ac diam ut lacus semper scelerisque quis at felis. Donec varius orci a dui fringilla vitae suscipit massa placerat. Aliquam interdum dui vel neque aliquam nec euismod eros placerat. Morbi id nibh erat. Class aptent taciti sociosqu ad litora torquent per conubia nostra, per inceptos himenaeos. In hac habitasse platea dictumst. Praesent velit augue, adipiscing dapibus faucibus in, accumsan a massa.
Curabitur ac sapien ante. Phasellus ac diam ut lacus semper scelerisque quis at felis. Donec varius orci a dui fringilla vitae suscipit massa placerat. Aliquam interdum dui vel neque aliquam nec euismod eros placerat. Morbi id nibh erat. Class aptent taciti sociosqu ad litora torquent per conubia nostra, per inceptos himenaeos. In hac habitasse platea dictumst. Praesent velit augue, adipiscing dapibus faucibus in, accumsan a massa.
\subsection{Lindt}
Mauris dignissim lectus id tortor tempus sed ultricies nibh dignissim. Praesent id eleifend risus. In hac habitasse platea dictumst. Morbi ac nibh urna, eu pharetra nunc. Lorem ipsum dolor sit amet, consectetur adipiscing elit. Curabitur egestas, odio vel scelerisque vulputate, nibh nibh scelerisque risus, sit amet congue tellus metus sed libero. Praesent fringilla justo at lectus auctor ac euismod sapien lobortis. Nullam sed convallis leo.
\subsection{Milka}


Curabitur ac sapien ante. Phasellus ac diam ut lacus semper scelerisque quis at felis. Donec varius orci a dui fringilla vitae suscipit massa placerat. Aliquam interdum dui vel neque aliquam nec euismod eros placerat. Morbi id nibh erat. Class aptent taciti sociosqu ad litora torquent per conubia nostra, per inceptos himenaeos. In hac habitasse platea dictumst. Praesent velit augue, adipiscing dapibus faucibus in, accumsan a massa.
\subsubsection{Loerrach}
Phasellus cursus, eros et scelerisque euismod, enim arcu rhoncus neque, luctus accumsan odio est nec neque. Quisque ultricies mi ultricies risus adipiscing sit amet suscipit odio aliquet. Curabitur aliquam euismod mauris. Ut imperdiet lacinia sapien, id condimentum ligula imperdiet eu. Vivamus mollis posuere quam, sed bibendum massa aliquet sed. Donec volutpat luctus mi at commodo. Vivamus faucibus elementum sem a posuere. Donec feugiat, elit in dapibus pulvinar, nibh mi faucibus ipsum, ac volutpat lectus justo nec tellus. Cras in suscipit erat. Vivamus quis dui sem. Nunc posuere ullamcorper tellus, sodales vestibulum massa eleifend quis. Vivamus quis sapien nec metus convallis vulputate quis ac ligula. 
Curabitur ac sapien ante. Phasellus ac diam ut lacus semper scelerisque quis at felis. Donec varius orci a dui fringilla vitae suscipit massa placerat. Aliquam interdum dui vel neque aliquam nec euismod eros placerat. Morbi id nibh erat. Class aptent taciti sociosqu ad litora torquent per conubia nostra, per inceptos himenaeos. In hac habitasse platea dictumst. Praesent velit augue, adipiscing dapibus faucibus in, accumsan a massa.

\section{Kaffee}
Curabitur ac sapien ante. Phasellus ac diam ut lacus semper scelerisque quis at felis. Donec varius orci a dui fringilla vitae suscipit massa placerat. Aliquam interdum dui vel neque aliquam nec euismod eros placerat. Morbi id nibh erat. Class aptent taciti sociosqu ad litora torquent per conubia nostra, per inceptos himenaeos. In hac habitasse platea dictumst. Praesent velit augue, adipiscing dapibus faucibus in, accumsan a massa.

\subsection{Java} Curabitur ac sapien ante. Phasellus ac diam ut lacus semper scelerisque quis at felis. Donec varius orci a dui fringilla vitae suscipit massa placerat. Aliquam interdum dui vel neque aliquam nec euismod eros placerat. Morbi id nibh erat. Class aptent taciti sociosqu ad litora torquent per conubia nostra, per inceptos himenaeos. In hac habitasse platea dictumst. Praesent velit augue, adipiscing dapibus faucibus in, accumsan a massa.
Curabitur ac sapien ante. Phasellus ac diam ut lacus semper scelerisque quis at felis. Donec varius orci a dui fringilla vitae suscipit massa placerat. Aliquam interdum dui vel neque aliquam nec euismod eros placerat. Morbi id nibh erat. Class aptent taciti sociosqu ad litora torquent per conubia nostra, per inceptos himenaeos. In hac habitasse platea dictumst. Praesent velit augue, adipiscing dapibus faucibus in, accumsan a massa.

	


    

    % Ausgabe Literaturverzeichnis
%    \newpage

\markboth{}{}%
\renewcommand{\baselinestretch}{.90}

%Quellen
\begin{thebibliography}{20}
%Reihenfolge: Personen, Internes  B"ucher, Magazine, Internetquellen
%%%%%%%%%   Personen 
\bibitem{myComp}
	Interne Dokumente, Marktanalysen, Strategiepläne und Gespräche mit Mitarbeitern\\
	aus dem eigenen Unternehmen
	<MONAT><jahr>, <Unternehmen Name, ORT>

%%%%%%%%%   Vorlesungsunterlagen
\bibitem{vl}
	Vorlesungsunterlagen <VL-Titel> \\
	Dozent: <Titel Name> \\
	<von> - <bis>,  DHBW-Lörrach

%%%%%%%%%   Buecher 
\bibitem{book}
	<autor> \\
	<Buch Titel> (1962)\\
	<Verlagsname>; Auflage: <wievielete Auflage><ErscheinungsDatum> \\
	ISBN-10: <ISBN Nummer>
	
\bibitem{knuth}
	Donald E. Knuth\\
	The Art of Computer Programming, Volumes 1-4\\
	Addison-Wesley Longman, Amsterdam; Auflage: 3rd edition. (3. März 2011)\\
	ISBN-10: 0321751043




%%%%%%%%%    Magazine
\bibitem{ix_110}
	ix - Magazin fuer Profesionelle Informationstechnik  \\
	1. Ausgabe Januar 2010, Seite 42\\
	Heise Zeitschriften Verlag GmbH \& Co. KG 

%%%%%%%%%    Internetquellen Paper
\bibitem{blogfefe}
	Blog von Felix von Leiter(fefe)\\
	Debian stolpert mal wieder über einen schlechten Zufallszahlengenerator.  \\
	Einsichtnahme: 30.04.2012  \\
	\url{http://blog.fefe.de/?ts=b1628b67}
   
 

\end{thebibliography}

      
    
    \bibliography{./literatur.bib}
    \clearpage
    
    % erzeuge Erklaerung    
    \makestatement

	
\end{document}




